%
\begin{abstract}
\addchaptertocentry{\abstractname} % Add the abstract to the table of contents 

Mass Spectrometry Imaging (MSI) is a revolutionary molecular microscopy technique that maps the spatial distribution of biomolecules within intact tissues, offering profound insights for precision medicine. However, the immense high-dimensional data generated by MSI is currently bottlenecked by heuristic preprocessing pipelines, specifically manual "peak picking." These classical algorithms introduce severe subjective bias, frequently discarding critical low-abundance disease biomarkers and distorting non-linear spectral manifolds. To overcome this limitation, this research proposes "Spatial-msiPL," a Spatially-Aware Variational Autoencoder designed for autonomous peak learning. By mathematically integrating physical neighborhood covariance into the encoding process via contextual input vectors and an explicit spatial coherence penalty, this lightweight 1D model aims to isolate true biomolecular signals from stochastic matrix noise. The proposed methodology will benchmark the extracted spatial ion images against expert-annotated structural segmentation masks using the mean Spatial F1-score (mSCF1). If successful, this framework will contribute a computationally efficient, spatially-aware alternative to heavy 3D convolutional networks, eliminating human bias and improving the clinical scalability of spatial omics analysis.
\end{abstract}